\section{Related Work}
\label{sec:related}

\subsection{Graph Neural Networks in Financial Crime \&  AML}
Graph Neural Networks (GNNs) have emerged as the leading computational paradigm for network-based fraud detection~\cite{weber2019anti,savage2016anomaly,dou2020enhancing}. Weber et al.~\cite{weber2019anti} pioneered the application of Graph Convolutional Networks (GCN) and GraphSAGE~\cite{hamilton2017inductive} to the Elliptic Bitcoin transaction dataset, proving that structural connectivity conveys strong predictive signals for illicit entity classification. However, early benchmarks predominantly operated on homogeneous, transaction-to-transaction or account-to-account graphs. In enterprise AML environments, transactions represent heterogeneous relationships between distinct real-world entities, including individual customer profiles, hardware identifiers, network endpoints, merchant points-of-sale, and physical banking infrastructure.

\subsection{Heterogeneous Graph Representation Learning}
To model diverse entity ecosystems, heterogeneous GNN architectures have been developed, including Relational Graph Convolutional Networks (R-GCN)~\cite{schlichtkrull2018modeling}, Heterogeneous Graph Attention Networks (HAN)~\cite{wang2019heterogeneous}, and Heterogeneous Graph Transformers (HGT)~\cite{hu2020heterogeneous}. HGT parameterizes type-specific projection matrices for query, key, value, and message computation, dynamically learning meta-relation attention weights without requiring predefined metapaths. While HGT has demonstrated remarkable representational capacity on academic citation and e-commerce graphs, its application to multi-bank financial networks under strict zero-leakage constraints remains under-explored.

\subsection{Target Leakage in Machine Learning Benchmarks}
Data leakage—the inclusion of information in training data that would not legitimately be available at inference time—represents a pervasive challenge in empirical machine learning~\cite{kaufman2012leakage,kapoor2023leakage}. In graph-based financial forensics, leakage manifests in two distinct modalities: (1) \textit{feature-level leakage}, where historical aggregation windows inadvertently compute statistics over future or ground-truth label targets (e.g., historical fraud counter rates), and (2) \textit{topological leakage}, where graph edges (such as explicit crime-ring memberships or SAR investigation case links) are constructed exclusively around known fraudulent entities. When evaluated on leaked graphs, models achieve near-perfect metrics that fail catastrophically upon real-world deployment. Establishing explicit zero-leakage evaluation protocols is therefore a critical prerequisite for credible financial forensics research.

\subsection{Continual Graph Learning \&  Catastrophic Forgetting}
Financial crime networks evolve dynamically to circumvent automated rule thresholds and deployed classification boundaries~\cite{wang2020streaming}. Updating neural architectures on non-stationary data streams frequently triggers \textit{catastrophic forgetting}~\cite{kirkpatrick2017overcoming,parisi2019continual,zenke2017continual}. Continual learning paradigms mitigate forgetting via architectural expansion, experience replay~\cite{zhang2022continual}, or regularization-based parameter consolidation. Elastic Weight Consolidation (EWC)~\cite{kirkpatrick2017overcoming} computes the diagonal of the empirical Fisher Information Matrix to identify and penalize deviations from parameters critical to previous tasks. This paper formalizes EWC within an attention-based heterogeneous graph framework for financial crime intelligence.
